\section{Annexes}

\subsection{Maquettes de référence}
\label{sec:maquettes}

Les maquettes ci-dessous illustrent les écrans décrits au §\ref{sec:fenetres}. Elles ne sont pas
normatives~: en cas de désaccord, ce sont les exigences du cahier des charges et la conception qui font foi.
Elles fournissent les valeurs de la table de constantes de chaque plateforme (M-04). Chaque plateforme est
présentée avec ses treize vues~: les onze écrans du cahier des charges, l'écran de liaison au compte
(\refdec{DEC-03}) et la liste des pièces hors connexion (\refexig{TPB-ACT-006}). La Pebble~2 Duo (flint) a
le rendu de la Pebble~2 (diorite).

\newcommand\maquette[2]{%
  \begin{figure}[p]
    \centering
    \includegraphics[width=\linewidth, height=0.86\textheight, keepaspectratio]{maquettes/#1.png}
    \caption{#2}
  \end{figure}}

\maquette{emery}{Pebble Time~2 (emery, 200\,×\,228, couleur)~: plateforme de référence}
\maquette{basalt}{Pebble Time et Time Steel (basalt, 144\,×\,168, couleur)}
\maquette{chalk}{Pebble Time Round (chalk, 180\,×\,180, rond, couleur)}
\maquette{gabbro}{Pebble Round~2 (gabbro, 260\,×\,260, rond, couleur)}
\maquette{diorite}{Pebble~2 et Pebble~2 Duo (diorite et flint, 144\,×\,168, noir et blanc)}

\clearpage
\subsection{Arborescence du projet}

\begin{verbatim}
tado-pebble/
  package.json            métadonnées, plateformes cibles, clés de messages
  src/c/
    main.c                M-01  cycle de vie
    comms.c               M-02  messagerie
    store.c               M-03  magasin, stockage persistant
    platform.c            M-04  table de constantes
    theme.c               M-05  couleurs, motifs
    list.c                M-06  moteur de liste (rect, rond)
    win_pieces.c ... win_liaison.c    M-07  une fenêtre par écran
    text.c                M-08  textes, formats
    actionbar.c           M-09
    banner.c              M-10
  src/pkjs/
    index.js              J-01  messagerie
    link.js               J-02  liaison
    tado.js               J-03  client Tado
    cache.js              J-04
    home.js history.js comfort.js alerts.js control.js   J-06 à J-10
    views.js              J-11
  resources/              images, ressources de textes fr et en
  test/                   tests du pkjs avec faux client Tado
\end{verbatim}

\subsection{Points ouverts}
\label{sec:points-ouverts}

\listepointsouverts
